% Affine registration -- version written with ChatGPT.
% Pulled in by sections/registration/affine-registration.tex.

\subsubsection{ChatGPT}          % second version of the affine-registration text
\label{sec:affine_chatgpt}       % \secref{sec:affine_chatgpt}

Following template construction, the airway segmentation volumes were spatially
aligned to a common anatomical reference using affine image registration. The
registration was performed using the Advanced Normalization Tools (ANTs)
framework through its Python interface, ANTsPy. ANTs is an open-source
framework for quantitative medical image analysis whose core functionality
includes image registration and template construction
\cite{avants2009ants,tustison2021antsx}. For each subject, the corresponding
three-dimensional airway segmentation volume from the AIIB23 dataset was
considered as the \emph{moving} image, whereas the previously constructed
airway template, \texttt{template\_22.nii.gz}, was used as the \emph{fixed}
image. In the implemented pipeline, all available input volumes were processed
independently and registered to this fixed reference.

% ------------------------------------------------------------
\paragraph{Affine transformation.}  % definition of the affine map

The affine registration models global geometric differences between an
individual airway tree and the reference template. In three-dimensional
space, an affine transformation can be written as
\begin{equation}
  \mathbf{x}' = \mathbf{A}\mathbf{x} + \mathbf{t},
  \label{eq:affine_transform}    % linear part A plus translation t
\end{equation}
where $\mathbf{x} \in \mathbb{R}^{3}$ denotes a point in the original image
space, $\mathbf{x}' \in \mathbb{R}^{3}$ its transformed position,
$\mathbf{A} \in \mathbb{R}^{3\times3}$ is the linear component of the
transformation, and $\mathbf{t} \in \mathbb{R}^{3}$ is the translation
vector. The matrix $\mathbf{A}$ can represent rotation, scaling, and shearing,
whereas $\mathbf{t}$ accounts for translation
\cite{itk_software_guide}. Consequently, a three-dimensional affine
transformation contains twelve parameters, corresponding to the nine elements
of the linear transformation matrix and the three translation parameters. The
ANTsPy documentation explicitly defines its \texttt{Affine} registration
option as an affine transformation with 12 parameters
\cite{antspy_registration}.

In contrast to deformable registration, the affine transformation is global:
the same transformation is applied throughout the image. It can therefore
correct subject-specific differences in overall position, orientation, scale,
and global shear, but it does not introduce spatially varying local
deformations. Consequently, the relative geometric configuration of the
airway tree is preserved while the tree is brought into a common global
coordinate system \cite{itk_software_guide,avants2009ants}.

% ------------------------------------------------------------
\paragraph{Registration as an optimization problem.}  % transform + metric + optimiser

Let $F$ denote the fixed template and $M$ the moving subject image. The
registration seeks the affine transformation parameters
$\boldsymbol{\theta}$ that optimize a similarity-based registration objective
between the fixed image and the transformed moving image. In general, the
registration problem can be expressed as
\begin{equation}
  \boldsymbol{\theta}^{*}
  = \underset{\boldsymbol{\theta}}{\operatorname{arg\,opt}}\;
  \mathcal{C}\left(F,\, M \circ T_{\boldsymbol{\theta}}\right),
  \label{eq:registration_optimization}  % optimise the objective over the affine parameters
\end{equation}
where $T_{\boldsymbol{\theta}}$ is the affine transformation parameterized by
$\boldsymbol{\theta}$, $M \circ T_{\boldsymbol{\theta}}$ denotes the moving
image after transformation, and $\mathcal{C}$ denotes the registration
objective. Image registration frameworks such as ITK formulate registration
through the combination of a transformation model, a similarity metric, and
an optimization procedure \cite{itk_software_guide,avants2014itkv4}.

In the present implementation, the registration is invoked through the ANTsPy
interface as

\begin{lstlisting}
ants.registration(
    fixed=fixed,
    moving=moving,
    type_of_transform='Affine'
)
\end{lstlisting}

\noindent where the fixed image is the image defining the reference space and
the moving image is the image mapped to that space \cite{antspy_registration}.
The choice \texttt{type\_of\_transform='Affine'} therefore specifies that the
estimated transformation belongs to the affine transformation family rather
than to a deformable model.

The implementation does not explicitly specify additional registration
parameters such as the similarity metric, optimization schedule, number of
iterations, or multiresolution settings. These are therefore determined by
the ANTsPy \texttt{Affine} registration preset associated with the software
version used in the experiment. The ANTsPy documentation exposes these
parameters explicitly, including the affine metric, sampling strategy,
iteration schedule, shrink factors, and smoothing sigmas
\cite{antspy_registration}. This distinction is important for reproducibility:
only parameters explicitly specified in the experimental code should be
reported as manually selected parameters.

% ------------------------------------------------------------
\paragraph{Mapping and resampling to template space.}  % from transform to output volume

After estimating the affine transformation, the transformed moving image is
resampled onto the spatial domain of the fixed image. In the ITK registration
framework, the moving image is evaluated at the physical coordinates induced
by the estimated transformation, and an interpolator is used when mapped
locations do not coincide with the original voxel grid
\cite{itk_software_guide}. The purpose of this resampling step is to obtain a
representation of the moving image on the grid of the fixed image, allowing
the registered volumes to share a common spatial reference.

In the implemented pipeline, the registered volume is obtained from the
\texttt{warpedmovout} output returned by \texttt{ants.registration}. According
to the ANTsPy documentation, \texttt{warpedmovout} is the moving image warped
into the space of the fixed image \cite{antspy_registration}. The resulting
volume is subsequently written as a new NIfTI file, thereby producing an airway
segmentation expressed in the coordinate system of the template.

The complete procedure can therefore be summarized as
\begin{equation}
  \begin{split}
    &\text{subject airway segmentation}
     \;\xrightarrow{\text{affine registration}}\;
     \text{estimated affine transformation} \\
    &\qquad\xrightarrow{\text{resampling}}\;
     \text{registered segmentation in template space}.
  \end{split}
  \label{eq:affine_pipeline}     % split over two lines so it fits the page width
\end{equation}

The resulting registered airway volumes share a common spatial reference and
can consequently be used for subsequent inter-subject analysis, including
morphological comparison, shape analysis, and representation learning.

% ------------------------------------------------------------
\paragraph{Implementation and computational setup.}  % parallel execution and outputs

The registration was applied independently to all available AIIB23 airway
segmentation volumes. The implementation uses a parallelized execution
strategy, with 12 registration jobs and 4 ITK threads assigned to each job.
These settings concern the computational execution of the pipeline and do not
change the mathematical definition of the affine transformation itself.

The registered volumes were saved in the specified output directory using the
subject identifier followed by the suffix \texttt{\_R.nii.gz}. These files
represent the final airway segmentations after affine alignment to the common
template and were used as inputs to the subsequent processing stages.
